% Generated draft: requires mathematical and source review.
\documentclass[11pt]{article}
\usepackage[margin=25mm]{geometry}
\usepackage{fontspec}
\setmainfont{Latin Modern Roman}
\usepackage{amsmath,amssymb,mathtools,unicode-math}
\setmathfont{Latin Modern Math}
\usepackage{enumitem,needspace,xurl,hyperref}
\setlength{\parindent}{0pt}
\setlength{\parskip}{5pt}
\setlength{\emergencystretch}{3em}
\newcommand{\N}{\mathbb{N}}
\newcommand{\Z}{\mathbb{Z}}
\newcommand{\Q}{\mathbb{Q}}
\newcommand{\R}{\mathbb{R}}
\newcommand{\C}{\mathbb{C}}
\newcommand{\F}{\mathbb{F}}
\newcommand{\A}{\mathbb{A}}
\newcommand{\PP}{\mathbb{P}}
\newcommand{\E}{\mathbb{E}}
\newcommand{\eps}{\varepsilon}
\newcommand{\dd}{\,\mathrm{d}}
\begin{document}
\section*{UnsolvedMath ID: 30001472}
\textbf{Source problem code:} catalogue entry
\subsection{Definitions and mathematical statement}
Let \(M\) be a finite-volume hyperbolic \(n\)-manifold, and let \(\mathcal C(M)\) denote its set of cusps. Represent each cusp \(C\) by a cusp neighborhood and put \(P_C=\pi_1(C)\). For a finite-index subgroup \(L\leq P_C\), let \(C_L\to C\) denote the connected covering corresponding to \(L\); without chosen basepoints, the covering is determined by the conjugacy class of \(L\) in \(P_C\).

The supplied question asks whether one can choose a sublattice \(\Lambda_C\subseteq P_C\) for every cusp simultaneously so that every subsequent choice of sublattices \(\Lambda'_C\subseteq\Lambda_C\) can be realized by a single finite covering of \(M\).

Under the interpretation that a sublattice of a cusp group is a finite-index free abelian subgroup of rank \(n-1\), and that subsequent sublattices also have finite index, a precise formulation is: does there exist a family \((\Lambda_C)_{C\in\mathcal C(M)}\) satisfying
\[
\Lambda_C\leq P_C,\qquad [P_C:\Lambda_C]<\infty,\qquad \Lambda_C\cong\mathbb Z^{n-1},
\]
such that, for every family \((\Lambda'_C)_{C\in\mathcal C(M)}\) satisfying
\[
\Lambda'_C\leq\Lambda_C,\qquad [\Lambda_C:\Lambda'_C]<\infty,
\]
there exists a finite covering \(p:M'\to M\) for which, for every \(C\in\mathcal C(M)\) and every connected component \(D\) of \(p^{-1}(C)\), the restricted covering \(p|_D:D\to C\) is isomorphic over \(C\) to \(C_{\Lambda'_C}\to C\)?

This formulation interprets the phrase ``whose cusps over \(C\) are exactly the covers determined by \(\Lambda'_C\)'' as prescribing the covering type of every cusp above \(C\), with no prescribed number of such cusps. The source does not explicitly settle this multiplicity interpretation, the finite-index meaning of ``sublattice,'' or whether \(M'\) must be connected. The formulation above does not impose connectedness on \(M'\). For a cusp with nontrivial flat holonomy, \(P_C\) need not itself be free abelian; the interpretation above therefore places the lattice inside \(P_C\), rather than assuming that \(P_C\) is a lattice. The source also does not explicitly specify the range of \(n\), completeness, connectedness, orientability, or boundary conventions for \(M\). These conventions require confirmation before treating this reading as an unqualified statement of the intended conjecture. If there are no cusps, the cusp conditions are vacuous and the identity covering satisfies them.
\subsection{Short English statement}
The question asks whether sufficiently deep initial sublattices can be chosen in all cusp groups so that any further finite-index choices occur simultaneously as the covering types of all cusps in some finite cover of the manifold. The source leaves several conventions unresolved, especially the meaning of sublattice for cusps with holonomy and the meaning of ``exactly.'' Its supplied dated assessment reports the conjecture as known in dimension three and open in general higher dimensions; this assessment has not been independently verified.
\subsection{Sources}
\begin{itemize}
\item Local UnsolvedMath registry: supplied complete source record, permanent ID \(30001472\), exact source problem code \texttt{catalogue entry}, title ``Cusp Covering Conjecture for Hyperbolic Manifolds,'' in the Oberwolfach Reports Open Problems collection. The record gives proposed year \(2010\) and supplies no named proposer.
\item Supplied report citation: \emph{Non-positive Curvature and Geometric Structures in Group Theory} (2010), Oberwolfach Reports, pp.\ 1165--1224, DOI \texttt{10.4171/owr/2010/20}. Source URL and supplied literature reference ``OWR 2010/20'': \url{https://doi.org/10.4171/owr/2010/20}. The record describes the collection as drawn from workshop reports and their open-problem discussions; it supplies no exact problem page within this report.
\item Supplied literature reference ``Behrstock--Neumann'': \url{https://www.math.columbia.edu/department/neumann/preprints/hypgraph_final.1.0.pdf}. The record supplies no further bibliographic details or theorem locator for this reference.
\item Literature assessment dated exactly \texttt{2026-08-21}: the supplied curation labels the problem open, reports that the cusp-covering conjecture is known in dimension three, and describes the general higher-dimensional assertion as remaining open and conditional on major residual-finiteness input. This is supplied dated curation, not an independently verified literature assessment. The background explicitly requests independent verification of time-sensitive status and cited claims; it does not identify the precise residual-finiteness hypothesis.
\item The registry labels the statement \texttt{corrected\_verified} and reports a manual semantic, typographic, and LaTeX review of its cleanup. These are supplied provenance claims, not independent verification of the formulation or mathematical status.
\end{itemize}
\end{document}
